\Needspace{9\baselineskip}
\originalchapter{通信复杂度、密码与安全边界}
\originalsection{通信复杂度模型}
Alice知道$x$, Bob知道$y$, 共同计算$f(x,y)$. 复杂度计交换比特或量子比特, 本地计算通常免费, 允许随机性、预共享纠缠与错误概率必须先指定. 否则不同模型的上界下界不可比较. 例如分布式搜索可通过双方可逆查询组成一个目标oracle, Grover把查询次数从线性降为平方根, 但每次查询的寄存器传输需要计$\log N$量级通信. 不能把查询数直接当通信比特数.

量子指纹把长字符串映成$O(\log n)$比特的非正交态, 相同字符串态相同, 不同字符串重叠受编码距离控制. 接收方SWAP test的不同结果概率为$[1-|\langle f_x|f_y\rangle|^2]/2$, 可重复判相等. 它适用于同时发送到裁判等指定模型, 不是无条件把所有两方通信都压缩到对数.

\originalsection{BB84与信息--扰动}
BB84用Z或X基的0、1态, Bob随机选基, 公布基后保留同基结果. 截获再发且随机猜基的Eve在保留位引入25\%错误: 一半猜错基, 错基后Bob同基读数又一半错. 这只是特定攻击的诊断, 不是一般安全证明.

在独立比特与理想设备的渐近协议中, CSS式纠缠纯化把bit错误校正与phase错误校正联系到密钥安全. 单向理想BB84的标准下界为$1-2h_2(Q)$, 对称误差假设下约$Q<11\%$有正率. 实际有限样本、设备缺陷、认证及信息泄露会改变可用密钥; 量子信道本身不替代经典消息认证. 本卷安全论证限课程理想模型, 不作为现实系统部署认证.

\originalsection{不可克隆与量子认证}
若幺正克隆两个可能态,$U\ket\psi\ket0=\ket\psi\ket\psi$, $U\ket\phi\ket0=\ket\phi\ket\phi$, 内积保持要求$s=s^2$, 仅$s=0$或1. 因而非正交未知态不可完美克隆. Eve可与环境做弱交互取得部分信息而非完整克隆, 安全不能只用这一定理一句话证明. 量子一次一密用均匀随机Pauli使任意单比特平均为$I/2$, 多比特需$2n$经典随机比特; 保密不自动提供篡改检测, 后者需要认证码及密钥管理.
